\documentclass[11pt]{article}
\usepackage[a4paper,margin=25mm]{geometry}
\usepackage{amsmath,amssymb,amsthm,booktabs,longtable,hyperref}
\setlength{\emergencystretch}{1em}
\newtheorem{theorem}{Theorem}\newtheorem{proposition}[theorem]{Proposition}
\newcommand{\N}{\operatorname{N}}\newcommand{\C}{\mathrm C}
\title{Finite subgroup atlases: proofs, certified tables, and source qualifications}
\author{L054 synthesis of acquired components}\date{}
\begin{document}\maketitle
This manuscript concerns Connor and Leemans' seven-page author edition
of \emph{An atlas of subgroup lattices of finite almost simple groups}
(arXiv version 2, 29 November 2013; PDF dated 2 December 2013), and
collates its eight-page 2015 journal version, pages 259--266.
The maximal descent and recursive naming constructions are the authors'
ideas. The arguments below and the acquired complete proofs make their
interfaces explicit and correct several ancillary formulations.
No novelty or new historical measurement is asserted.

\paragraph{Main results and the atlas boundary.}
We recover the maximal-descent and structural-certificate constructions under
explicit complete local subgroup and exact-operation interfaces. The acquired
classified roster contains all 127 almost simple isomorphism classes of order
at most one million. Complete numerical tables and structural certificates
have been checked for every one of these groups, including 47 groups absent
from the current small companion roster. Seven larger applications are also
complete. Their sizes are recorded below.
\begin{center}
\begin{tabular}{lrr}\toprule
Domain & Groups & Subgroup-class rows\\\midrule
All almost simple groups of order at most $10^6$ & 127 & 10,694\\
Seven specified larger groups & 7 & 5,294\\\midrule
Complete reconstruction & 134 & 15,988\\\bottomrule
\end{tabular}
\end{center}
These are row occurrences in actual tables, not pairwise different abstract
subgroups. Of these rows, 13,520 belong to 87 current companion tables;
one printed label in the $A_{11}$ table is refuted and replaced by its
verified structure. The remaining 2,468 rows are additional complete
applications with generated structural descriptions.

This proves the small-order claim for a domain containing the source's
2012 online-ATLAS domain and supplies more than one hundred complete tables.
It is \emph{not} a reproduction or verification of the entire named historical
or current atlas. The original articles give no exact filename roster.
The current inventory of 116 companion tables, specifically attributed
outputs such as the O'Nan table, and unlinked numerical claims remain in a
separate ledger. No unpaid current table acquires its structural labels
through this main-result reconstruction. Historical file identity and the
reported historical timings are not asserted.

\section{The maximal-descent and weight argument}
All groups here are finite. For $H\leq G$, write $[H]_G$ for its conjugacy
class, ordered by containment of representatives. Suppose each processed
$H$ comes with an actual isomorphism $\phi_H:H\to\widetilde H$, a complete
cover of its maximal-subgroup conjugacy classes, and exact ambient conjugacy
tests. Lists may repeat a local class; duplicates must first be identified
by actual local conjugacy. Completeness is an explicit mathematical/software
input. It does not follow from a plausible list length or matching orders.

\begin{theorem}
Maintain a permanent registry of both processed and pending classes,
initially containing $G$. Process each registered parent once; pull back its
complete local maximal cover, fuse the resulting subgroups under conjugacy
in $G$, and register each unseen class. This terminates with exactly all
subgroup classes of $G$. Reachability among its actual maximal-inclusion
edges is the entire class containment relation.
\end{theorem}
\begin{proof}
An isomorphism and its inverse preserve subgroup containment, maximality,
conjugacy and normalizers. Hence every new class and edge is genuine.
There are finitely many classes, and no registered class is processed twice.
For completeness, refine any $K\leq H$ to a maximal subgroup chain.
At each step a member of the complete local cover is conjugate inside the
parent to the next chain term. Its ambient fusion only changes the
representative by conjugacy. Induction down this chain, starting with $G$,
discovers $[K]_G$. The same argument inside any larger representative proves
that containment implies reachability. The reverse implication follows
from transitivity of containment on classes.
\end{proof}
The permanent registry repairs the source prose's test against only the
current queue unless an additional scheduling invariant is supplied.
Already completed parents can be reused: their complete transported local
lists and exact fusion remain valid, while every unprocessed parent is
still an obligation. The acquired literal-state and continuation interfaces
implement this invariant. Named display strings alone are not embedded
subgroup data; a failed symmetric-group checkpoint demonstrated why actual
generators must be serialized.

For representatives $H,K$, let $M_1,\ldots,M_t$ be the distinct local
$H$-classes of maximal subgroups that fuse to $[K]_G$. Put
\[
 c_H=[G:\N_G(H)],\qquad
 a_{HK}=\sum_{j=1}^t[H:\N_H(M_j)],\qquad
 b_{KH}=c_Ha_{HK}/c_K.
\]
Orbit--stabilizer proves that $a_{HK}$ counts the actual maximal embeddings
inside $H$. Counting the pairs of conjugates $(H',K')$ with $K'$ maximal
in $H'$ in both orders proves the expression for $b_{KH}$ and its integrality.
Local classes that fuse must be added, rather than counted once.
Duplicate local representatives must not be added twice. The exact bounds
$a_{HK}\leq c_K$ and $b_{KH}\leq c_H$ follow immediately.

The class-poset covers are obtained by transitive reduction of the complete
maximal-edge graph. On a cover, every embedding is maximal, so these weights
are exactly the source's $n_{HK},n_{KH}$. Outside covers, the displayed
maximal-edge weights do not automatically count every embedding.
The acquired robust local-list proof and retained-inventory reweighting
pay these distinctions. In particular, an early duplicate local-class
output for the $\operatorname{P\Gamma U}_3(5)$ model inflated one computed
weight from $(3,1)$ to $(6,2)$; the corrected complete replay matches the
source. This was our execution issue, not a source-table error.

\section{Recursive names as exact existence assertions}
A structural expression is interpreted recursively. Parentheses in the
native source are retained; otherwise products associate to the left.
A cross denotes a direct product, a colon a split extension, and the
centered dot a nonsplit extension. The source also uses a different full
stop for an extension whose splitting has not yet been specified.
Names without the action or factor data need not determine a unique group.

For $N\lhd H$, $Q=H/N$, choose a section $s:Q\to H$ with $s(1)=1$.
Set $\alpha_q(n)=s(q)ns(q)^{-1}$ and
$f(q,r)=s(q)s(r)s(qr)^{-1}\in N$. Then $N\times Q\to H$,
$(n,q)\mapsto ns(q)$, is bijective, and multiplication is
\[
 (n,q)(m,r)=(n\alpha_q(m)f(q,r),qr).
\]
The acquired complete extension proof verifies associativity exactly from
$\alpha_q\alpha_r=\operatorname{Inn}(f(q,r))\alpha_{qr}$ and
$f(q,r)f(qr,t)=\alpha_q(f(r,t))f(q,rt)$, with the identity normalizations.
An explicit complement proves splitting; a commuting complement proves a
direct product. Failure of a search to find a complement proves neither
nonsplitting nor absence of a specified structure.

At each nonterminal naming step choose a proper nontrivial normal subgroup.
Both child orders are smaller, so recursion terminates. Exact leaf
identifications and the verified extension data reconstruct the parent by
induction. Without those data the expression records only an existence
assertion. The source's preferences for short or recognizable factors may
select one such description, but are not a canonical-isomorphism theorem.

Every accepted row uses an actual embedded subgroup. An earlier named
model is transported by an explicit bijection fixing identity and respecting
right multiplication by every ordered generator. Induction on words proves
that this is a homomorphism; bijectivity proves it is an isomorphism.
Complete Cayley-word maps, faithful smaller actions, elementary normal forms,
actual normal factors and complements, and exhaustive quotient-lift
obstructions supply the concrete certificates. Each criterion has its
acquired complete proof. Source typography is checked separately: a flattened
superscript cannot silently become a cyclic order, and a dot cannot silently
be replaced by a colon. Repeated descriptions may use different concrete
models, each with its own exact transport.

\section{Finite application and retained inputs}
The attached current small-table roster has 80 tables and 8,226 subgroup-class
rows. Each numerical route fixes an actual named group, a complete subgroup
inventory, actual embedded generators, local class fusion, both incidence
weights, and a bijection to the native source rows. Each structural route
checks the source expression and every required normal factor, quotient,
complement or nonsplitting obstruction. Their complete proof texts and
separate author and reviewer records are acquired without another census.

\begin{theorem}
Subject to the acquired classification of finite simple groups, their exact
orders and outer automorphism groups, there are 127 almost simple
isomorphism classes of order at most one million. For each class, the
acquired applications give its full subgroup-class inventory, maximal
incidences and a certified structural description of every subgroup class.
Together these tables contain 10,694 subgroup-class rows.
\end{theorem}
\begin{proof}
The acquired small-order roster proof first bounds the simple socle using
the exact family orders and the retained classified exceptional list.
There are 56 socles in this range. For a fixed socle $S$, an almost simple
group is the full inverse image of a subgroup of
$\operatorname{Out}(S)=\operatorname{Aut}(S)/S$. An isomorphism preserves
the characteristic simple socle and therefore identifies such groups
exactly by the corresponding outer-subgroup conjugacy relation. The complete
small outer-group calculations give 127 cases. The complete roster proof,
classification locators and arithmetic checks are included in the proof
collection; no new proof of the classification is claimed here.

The 80 current small-table applications cover their exact named targets.
The remaining 47 cases have 28 complete parent descents and 19 exact
index-two restrictions from those parent inventories. The descent theorem
above and the acquired restriction theorem prove completeness. Explicit
named-model identities join each group to its roster position. Complete
generator-edge transports, normal factors, complements and nonsplitting
obstructions prove the 2,468 additional structural rows. Thus all 127
roster positions have complete numerical and structural applications.
\end{proof}

\begin{proposition}
The following seven larger groups have complete numerical tables and
structural certificates, with the stated correction for $A_{11}$.
\end{proposition}
\begin{center}
\begin{tabular}{lr}\toprule
Group & Subgroup-class rows\\\midrule
$A_{10}$ & 430\\
$A_{11}$ & 788\\
$M_{23}$ & 204\\
$\operatorname{PSp}(4,5)$ & 307\\
$\operatorname{PSp}(4,7)$ & 603\\
$\operatorname{Sp}(6,2)$ & 1,369\\
$S_{10}$ & 1,593\\\bottomrule
\end{tabular}
\end{center}
\begin{proof}
Each application acquires a complete numerical inventory and an exact
bijection to the companion rows, then verifies every displayed structural
assertion with the literal subgroup generators. The acquired mathematical
proofs justify the reused transport, support, extension and finite
normal-form tests. The seven independent table reviews bind all 5,294
rows and their native expressions. The $A_{11}$ review recovers 787 printed
labels and proves the remaining printed label false, supplying its actual
structure instead. The exact correction is stated below. No inference from
an unverified group name or from matching orders substitutes for these
certificates.
\end{proof}

Only $M_{23}$ is a sporadic group among these seven larger applications.
We do not infer coverage of the source's plural historical larger-sporadic
examples from that one case. The quantified main-result reconstruction and
the correctness of every individually named historical output are separate
claims. The latter is not made here.

The source's Section 4 reports all groups in its ATLAS-restricted small-order
domain, some larger groups, and more than one hundred complete outputs.
The 127-group theorem proves a stronger small-order mathematical statement;
the seven further applications supply explicit larger cases. No historical
membership assumption is needed for that quantified conclusion. Conversely,
these conclusions do not validate every current companion PDF. The retained
index has 133 group rows, 130 linked occurrences and 116 distinct primary
links, with a printed update date of 4 August 2013. Its documentary date
does not establish historical byte identity. The O'Nan landing page
specifically attributes its current linked table to this article's programs;
that particular companion obligation remains recorded and unpaid. The
source-claim map preserves all such boundaries. The companion index bounds
the simple socle order $|S|$, whereas the article bounds $|G|$; our
small-order theorem does not establish that stronger companion domain.
The unlinked
Suzuki and Fischer numerical assertions are retained separately. The completed $S_{11}$ numerical
inventory remains additional evidence with its structural labels unproved.

Named constructors, complete external local covers or previously accepted
ambient inventories, and the exact group-operation interfaces remain
explicit premises. Source counts are not used as completeness filters.
All earlier mathematical authorship and distinct reviews are retained;
this synthesis neither repeats the prior proofs nor represents itself as
an independent review of its own earlier components.

\section{Exact source corrections and attribution}
\paragraph{Class posets and joins.}
The subgroup join is $\langle H,K\rangle$, not ordinary set union.
The conjugacy-class poset need not be a lattice, contrary to the literal
sentence on author page 2 and journal page 261.
For an elementary counterexample take $G=A_5$. The classes of subgroups
$C_2$ and $C_3$ have both $S_3$ and $A_4$ as upper bounds: embed $S_3$
on three letters and compensate every odd permutation by swapping the
other two; embed $A_4$ on four letters. The relevant $C_2$ subgroups are
conjugate in $A_5$, as are the $C_3$ subgroups. Indeed an odd conjugator
may be made even using an odd permutation centralizing the relevant element.
Any least upper bound would lie below both $S_3$ and $A_4$, and its order
would be divisible by six and divide six. It would therefore have order six.
But $A_4$ has no subgroup of order six: such a subgroup would have index two,
whereas every homomorphism $A_4\to C_2$ kills the generating 3-cycles.
Thus no such least upper bound exists. The accepted literal $A_5$ catalogue
provides the same incidence obstruction.

\paragraph{Composition and deterministic names.}
Strict normal inclusions alone do not define a composition series: the
factors must be simple, equivalently the lower term is maximal normal in
the next one. Jordan--H\"older gives the multiset of simple factors up to
isomorphism, not uniqueness of the intermediate subgroups. For example,
$C_2^2$ has three different order-two intermediate terms with the same two
composition factors. These are qualifications to author page 4 and journal
page 263, not new classification results.

The impossibility sentence about deterministic names on those pages is
also too strong literally. Number the identity by zero, enumerate all
remaining labelings of a finite group's complete multiplication table,
and take the lexicographically least resulting table as its name.
Isomorphic groups yield the same collection of relabeled tables; equality
of the minima supplies an isomorphism. This finite deterministic construction
is not an efficient or useful short notation. It shows that ambiguity of
factor descriptions is not undecidability of finite group identification.
The source's example of nonisomorphic groups with the same normal-section
information remains attributed to its earlier cited paper.

\paragraph{The group-name chain.}
Both native author page 3 and journal page 262 print a chain beginning
$S_4(3)\cong U_4(2)$ and including $O^-_6(3)$ and $W(E_6)$.
As printed the chain is false. The independently acquired named
$\operatorname{PSU}(4,2)$ certificate has order $25920$, whereas the acquired
$E_6$ Weyl calculation has order $27\cdot1920=51840$.
These orders preclude the asserted isomorphism with $W(E_6)$.
We do not silently replace the field parameter or a full Weyl group by
an index-two subgroup. The valid general point---that one group may have
several geometric descriptions---does not require that faulty chain.

\paragraph{The $A_{11}$ companion correction.}
Row 186 prints the left-associated expression $(2.A_5):2$.
Its order-120 normal section would be the nonsplit double cover of $A_5$.
The independently reviewed application proves that this section cannot
embed in $A_{11}$ under the stated convention, and checks the actual
order-240 subgroup by a complete literal certificate. Its structure is
$(A_5\times C_2)\mathbin{\cdot}C_2$, equivalently
$C_2\mathbin{\cdot}(A_5:C_2)$.
The printed expression is retained as a refuted source assertion; it is
not silently reassociated. This correction concerns a companion label,
not a newly solved open conjecture.

\paragraph{Source examples and historical reports.}
The examples $S_5=A_5:C_2$ and $A_4=C_2^2:C_3$ have their standard actual
normal factors and complements, also present in the acquired registry.
The groups $S_5$ and $A_5\times C_2$ have the same two composition
factors but centers of orders one and two, respectively, so are not
isomorphic. A coordinate flag of $C_2^5$ gives the displayed product of
five order-two factors, each split and commuting.
The label $5\times2\mathbin{\cdot}2:2$ can describe both $D_{40}$ and
$C_5\times D_8$: use their cyclic order-20 normal subgroup, its order-10
subgroup, and an order-two complement. The intermediate index-two
extension is nonsplit. The final action in the first case inverts the
order-five factor and in the second case centralizes it. The centers have
orders two and ten, respectively, so the groups are not isomorphic.
This pays the source's concrete illustration of nonunique notation.

The introduction's earlier sporadic subgroup computations, the small-group
count near order 1024, software-database descriptions and historical
normal-section example remain attributed antecedents. The concrete
$L_3(7):2$ degree-reduction and machine-time/memory report, and the journal's
qualitative performance comparison, are empirical source reports.
Correctness of the algorithm does not prove a universal speed or memory
advantage, and no historical timing is reproduced or credited here.
The source-clause register lists every such boundary separately.

\clearpage
\section*{Exact current small-table roster}
\small
\begin{longtable}{lrr}
\toprule Current file locator & Group order & Rows\\\midrule\endfirsthead
\toprule Current file locator & Group order & Rows\\\midrule\endhead
\texttt{alt5} & 60 & 9\\
\texttt{sym5} & 120 & 19\\
\texttt{psl27} & 168 & 15\\
\texttt{pgl27} & 336 & 23\\
\texttt{alt6} & 360 & 22\\
\texttt{psl28} & 504 & 12\\
\texttt{psl211} & 660 & 16\\
\texttt{m10} & 720 & 25\\
\texttt{pgl29} & 720 & 26\\
\texttt{psigmal29} & 720 & 56\\
\texttt{psl213} & 1092 & 16\\
\texttt{pgl211} & 1320 & 29\\
\texttt{pgammal29} & 1440 & 65\\
\texttt{psigmal28} & 1512 & 25\\
\texttt{pgl213} & 2184 & 30\\
\texttt{psl217} & 2448 & 22\\
\texttt{alt7} & 2520 & 40\\
\texttt{psl219} & 3420 & 19\\
\texttt{psl216} & 4080 & 21\\
\texttt{pgl217} & 4896 & 32\\
\texttt{sym7} & 5040 & 96\\
\texttt{psl33} & 5616 & 51\\
\texttt{psu33} & 6048 & 36\\
\texttt{psl223} & 6072 & 23\\
\texttt{pgl219} & 6840 & 36\\
\texttt{psl225} & 7800 & 37\\
\texttt{m11} & 7920 & 39\\
\texttt{psl216d2} & 8160 & 47\\
\texttt{psl227} & 9828 & 16\\
\texttt{psl33d2} & 11232 & 78\\
\texttt{pgammau33} & 12096 & 100\\
\texttt{pgl223} & 12144 & 34\\
\texttt{psl229} & 12180 & 22\\
\texttt{psl231} & 14880 & 29\\
\texttt{pgl225} & 15600 & 43\\
\texttt{psigmal225} & 15600 & 103\\
\texttt{psl225d2} & 15600 & 39\\
\texttt{pgammal216} & 16320 & 69\\
\texttt{pgl227} & 19656 & 32\\
\texttt{alt8} & 20160 & 137\\
\texttt{psl34} & 20160 & 95\\
\texttt{psu42} & 25920 & 116\\
\texttt{suz8} & 29120 & 22\\
\texttt{psigmal227} & 29484 & 38\\
\texttt{pgl231} & 29760 & 44\\
\texttt{pgammal225} & 31200 & 115\\
\texttt{psl232} & 32736 & 24\\
\texttt{psigmal34} & 40320 & 152\\
\texttt{psl34d21} & 40320 & 151\\
\texttt{psl34d23} & 40320 & 102\\
\texttt{sym8} & 40320 & 296\\
\texttt{pgammau42} & 51840 & 350\\
\texttt{pgammal227} & 58968 & 67\\
\texttt{pgl34} & 60480 & 100\\
\texttt{psu34} & 62400 & 34\\
\texttt{psl34d2e2} & 80640 & 288\\
\texttt{suz8d3} & 87360 & 39\\
\texttt{m12} & 95040 & 147\\
\texttt{pgammal34} & 120960 & 226\\
\texttt{psl34d6} & 120960 & 143\\
\texttt{psl34ds3} & 120960 & 142\\
\texttt{psu34d2} & 124800 & 80\\
\texttt{psu35} & 126000 & 80\\
\texttt{pgammal232} & 163680 & 30\\
\texttt{j1} & 175560 & 40\\
\texttt{alt9} & 181440 & 223\\
\texttt{m12d2} & 190080 & 213\\
\texttt{psl34dD12} & 241920 & 409\\
\texttt{psigmau34} & 249600 & 120\\
\texttt{psigmau35} & 252000 & 148\\
\texttt{psl264} & 262080 & 76\\
\texttt{sym9} & 362880 & 554\\
\texttt{psl35} & 372000 & 140\\
\texttt{pgu35} & 378000 & 93\\
\texttt{m22} & 443520 & 156\\
\texttt{j2} & 604800 & 146\\
\texttt{psl35d2} & 744000 & 208\\
\texttt{pgammau35} & 756000 & 244\\
\texttt{m22d2} & 887040 & 490\\
\texttt{sp44} & 979200 & 496\\
\midrule Total: 80 tables &&8226\\\bottomrule
\end{longtable}

\normalsize
\end{document}
